% =====================================================================================================
\section{Overview of the workflow}\label{sec:overview}
% =====================================================================================================

\subsection{The task}
Each X-ray binary in the project is labelled from external evidence:
\begin{itemize}
\item \emph{black hole} (BH): a dynamical mass measurement;
\item \emph{neutron star} (NS): type-I X-ray bursts or pulsations.
\end{itemize}
The workflow turns single X-ray observations of these sources into feature vectors, using:
\begin{itemize}
\item RXTE/PCA spectra and light curves;
\item NICER/XTI event files;
\item MAXI/GSC daily light curves.
\end{itemize}
Classifiers are trained on the feature vectors. Their scores are always computed for sources that the model did not see
in training, and they are evaluated with statistics that treat the \emph{source}, not the observation, as the independent
unit. Everything is driven by Python scripts. HEASoft is used only where no Python equivalent exists: the XSPEC
cross-check, the HI4PI $N_\mathrm{H}$ look-up and the NICER 3C50 background.

\subsection{The workflow at a glance}
\Cref{fig:workflow} shows the whole chain, and \cref{tab:stages} lists the stages with their scripts and outputs. The
chain has five parts, each described in later sections:
\begin{enumerate}
\item \emph{External databases} (\cref{sec:acquisition}): source catalogues, observation catalogues and data archives.
\item \emph{Sample construction} (\cref{sec:samples}): which sources, and which of their observations.
\item \emph{Processing and feature extraction} (\cref{sec:rxte,sec:nicer}).
\item \emph{Data model} (\cref{sec:datamodel}): stored tables and the matrices given to the models.
\item \emph{Modelling and evaluation} (\cref{sec:ml,sec:evaluation}).
\end{enumerate}
\Cref{sec:operations} covers running, logging and reproduction.

\begin{figure}[htbp]
\centering
\begin{tikzpicture}[x=1cm, y=1cm,
  db/.style={fc data, text width=2.75cm, minimum height=17mm},
  st/.style={fc proc, text width=2.75cm, minimum height=17mm},
  md/.style={fc model, text width=2.75cm, minimum height=17mm},
  outb/.style={fc, draw=inkgrey, fill=boxgrey, text width=2.75cm, minimum height=17mm},
  lane/.style={font=\scriptsize\bfseries, text=inkgrey, anchor=east}]
  % row 1: external databases
  \node[db] (lab) at (1.55,0)  {\textbf{label catalogues}\\ BlackCAT, Marcel+2026, Galloway+2008, MINBAR, pulsar lists};
  \node[db] (pos) at (4.75,0)  {\textbf{positions}\\ CDS Sesame (SIMBAD); MINBAR source table};
  \node[db] (cat) at (7.95,0)  {\textbf{observation catalogues}\\ RXTE MissionLongData; TAP \texttt{xtemaster}, \texttt{nicermastr}};
  \node[db] (arc) at (11.15,0) {\textbf{data archives}\\ HEASARC RXTE tree; NICER AWS mirror; MAXI; GitHub};
  \node[db] (cal) at (14.35,0) {\textbf{calibration}\\ RXTE responses; NICER CALDB (3C50); HI4PI};
  % row 2: selection, retrieval, processing
  \node[st] (sel) at (1.55,-2.55)  {\textbf{select sources}\\ label rule, position match, minimum number of eligible observations};
  \node[st] (rank) at (4.75,-2.55) {\textbf{rank observations}\\ time-stratified groups, seeded order (Alg.~\ref{alg:sampling})};
  \node[st] (get) at (7.95,-2.55)  {\textbf{retrieve}\\ HTTPS download, range requests or streaming; SHA256 log};
  \node[st] (proc) at (11.15,-2.55) {\textbf{process}\\ net spectra, cospectra, colours, rms, $\nu_c$; background; bursts};
  \node[fc gate, text width=2.75cm, minimum height=17mm] (qual) at (14.35,-2.55) {\textbf{quality rules}\\ fail: try next candidate in the group};
  % row 3: data model, models, evaluation
  \node[outb] (tab) at (1.55,-5.1)  {\textbf{feature tables}\\ \file{features.npz}, \file{*_features.csv}, \file{observations.csv}};
  \node[md] (mat) at (4.75,-5.1)  {\textbf{design matrices}\\ $(\Xm,\yv,\gv)$; joins by ObsID; missing values};
  \node[md] (fit) at (7.95,-5.1)  {\textbf{models}\\ LR, RF (+ 6 others in the search); LOSO, nested or frozen};
  \node[md] (ev) at (11.15,-5.1)  {\textbf{evaluation}\\ AUC, balanced accuracy; source bootstrap; permutation tests};
  \node[outb] (res) at (14.35,-5.1) {\textbf{outputs}\\ \file{results/}, \file{figures/}, reports, decision log};
  \foreach \a/\b in {sel/rank, rank/get, get/proc, proc/qual, tab/mat, mat/fit, fit/ev, ev/res} \draw[fc arr] (\a) -- (\b);
  \draw[fc arr] (lab) -- (sel); \draw[fc arr] (pos.south) -- ++(0,-0.4) -| ([xshift=4mm]sel.north);
  \draw[fc arr] (cat) -- (rank); \draw[fc arr] (arc) -- (get); \draw[fc arr] (cal) -- (proc);
  \draw[fc arr] (qual.south) -- ++(0,-0.45) -| (tab.north);
  % side bar: configuration and provenance
  \node[fc, draw=tagprim, dashed, fill=white, text width=15.5cm, minimum height=8mm, anchor=north west] (cfg) at (0.18,-6.25)
    {\textbf{configuration and provenance (all stages):} \file{scripts/config.py} blocks selected by \code{XRB_VERSION};
     time-stamped plan \file{reports/preregistration_vN.md} before data; \file{logs/downloads.jsonl} (URL, bytes, SHA256);
     \file{reports/decision_log.md} (every change, with whether it was made after results); run scripts \file{run_vN.sh}.};
\end{tikzpicture}
\caption{End-to-end workflow. Amber: external databases; green: selection, retrieval and processing; blue: data model,
models and evaluation; grey: stored outputs; chamfered: a gate that can reject an observation. An observation that fails
the quality rules is replaced by the next ranked candidate of the same time group, so the sample is defined before any
spectrum or light curve is examined.}\label{fig:workflow}
\end{figure}

{\small\setlength{\tabcolsep}{3.5pt}
\begin{longtable}{@{}>{\raggedright\arraybackslash}p{2.6cm}>{\raggedright\arraybackslash}p{6.2cm}>{\raggedright\arraybackslash}p{3.6cm}>{\raggedright\arraybackslash}p{3.4cm}@{}}
\caption{Workflow stages, the scripts that implement them and their main outputs. Script numbers are those in
\file{scripts/}; RXTE and NICER have parallel chains.}\label{tab:stages}\\
\toprule
Stage & What happens & Scripts & Main outputs \\
\midrule\endfirsthead
\toprule Stage & What happens & Scripts & Main outputs \\ \midrule\endhead
\bottomrule\endfoot
Archive inspection & check the RXTE archive layout and file formats on two spectra & \file{01}, \file{05} & \file{reports/step2_fits_check.md} \\
Source tables & label lists to positions; match to observation catalogues & \file{03}, \file{13}, \file{20}, \file{28}, \file{35}, \file{38}, \file{41}, \file{47} & \file{data/*/sources.csv} \\
Candidate ranking & eligible observations, time-stratified groups, seeded order & \file{03}, \file{23}, \file{41}, \file{47}, \file{55}, \file{60} & \file{data/*/observation_candidates.csv} \\
Retrieval + processing (RXTE) & Standard Products, Standard-1, HEXTE, event mode; quality rules & \file{04}, \file{21}, \file{24}, \file{19}, \file{26}, \file{29}, \file{42}, \file{43} & \file{data/*/processed/features.npz}, timing tables \\
Retrieval + processing (NICER) & TAP, prefix/tail/streamed event files, features, 3C50 & \file{48}, \file{52}, \file{56}, \file{58}, \file{61}, \file{62} & \file{data/v9}--\file{v13/processed/*.csv} \\
Retrieval + processing (MAXI) & label freeze, Sesame matching, 1-day light curves, cuts; $N_\mathrm{H}$ & \file{37}, \file{38}, \file{45} & \file{data/v7c/maxi_points.csv}, \file{data/v8d/} \\
Bursts and states & Standard-1 burst detector, MINBAR match, rms state classes & \file{17}, \file{31}--\file{34}, \file{42} & \file{results/v4/burst_*}, \file{results/v7a/states.csv} \\
Data checks & integrity, response variation, XSPEC cross-check & \file{06}, \file{10} & \file{results/data_checks.md} \\
Modelling & LOSO, nested selection, frozen models & \file{07}, \file{11}, \file{16}, \file{27}, \file{30}, \file{44}, \file{49}, \file{54}, \file{57}, \file{59}, \file{63} & OOF prediction tables \\
Inference & source bootstrap, permutation tests, conformal sets & \file{09} and analysis scripts; \file{v4lib}--\file{v10lib} & \file{results/*/*bootstrap*.csv}, \file{*_nuc.csv} \\
Summaries & cross-round figures and tables & \file{15}, \file{40}, \file{46}, \file{51}, \file{64} & \file{figures/final/}, \file{results/final_*} \\
\end{longtable}}

\subsection{Analysis rounds and how a round runs}\label{sec:rounds}
The project grew in numbered rounds (v1--v14). \Cref{tab:rounds} lists them briefly. Every round from v2 on follows the
same procedure (\cref{fig:lifecycle}).

\begin{figure}[htbp]
\centering
\begin{tikzpicture}[x=1cm, y=1cm,
  st/.style={fc proc, text width=2.75cm, minimum height=15mm},
  gt/.style={fc gate, text width=2.75cm, minimum height=15mm}]
  \node[fc, draw=inkgrey, fill=white, text width=2.75cm, minimum height=15mm] (pl) at (1.55,0) {\textbf{plan}: \file{preregistration_vN.md}, written after all earlier results};
  \node[st] (cf) at (4.75,0) {\textbf{config block} in \file{config.py}, active only for \code{XRB_VERSION=vN}};
  \node[gt] (ck) at (7.95,0) {\textbf{earlier settings unchanged} (all earlier versions' values compared)};
  \node[gt] (ut) at (11.15,0) {\textbf{unit tests} \file{test_vNlib.py} on synthetic data (gates)};
  \node[fc data, text width=2.75cm, minimum height=15mm] (dl) at (14.35,0) {\textbf{selection and download} (logged)};
  \node[st] (fe) at (14.35,-2.45) {\textbf{features} with the frozen code of the round};
  \node[gt] (ic) at (11.15,-2.45) {\textbf{implementation checks} (IC): stop and report if one fails};
  \node[fc model, text width=2.75cm, minimum height=15mm] (an) at (7.95,-2.45) {\textbf{models and statistics} as pre-specified (primary, descriptive)};
  \node[fc, draw=inkgrey, fill=boxgrey, text width=2.75cm, minimum height=15mm] (rp) at (4.75,-2.45) {\textbf{results} CSV, figures, \file{report_vN_zh-TW.md}};
  \node[fc, draw=inkgrey, fill=boxgrey, text width=2.75cm, minimum height=15mm] (mg) at (1.55,-2.45) {\textbf{branch} \file{vN-analysis} merged to master (pull request)};
  \foreach \a/\b in {pl/cf, cf/ck, ck/ut, ut/dl, dl/fe, fe/ic, ic/an, an/rp, rp/mg} \draw[fc arr] (\a) -- (\b);
  \node[fc, draw=tagpost, dashed, fill=white, text width=15.4cm, minimum height=7mm, anchor=north west] at (0.17,-3.65)
    {\textbf{decision log} (\file{reports/decision_log.md}): every deviation, bug fix, added comparison or choice is recorded with
     the time and whether it was made after the round's results were seen; post hoc additions are labelled as such in reports.};
\end{tikzpicture}
\caption{Life cycle of one analysis round. The plan and the configuration are committed before the round's data are
downloaded. The implementation checks and unit tests are written before real results exist; if a check fails, the round
stops and the failure is reported rather than the rule being changed.}\label{fig:lifecycle}
\end{figure}

\begin{table}[htbp]
\centering\footnotesize
\caption{Analysis rounds (repository versions). ``Data'' is what the round added; plans are
\file{reports/preregistration_vN.md}.}\label{tab:rounds}
\begin{tabularx}{\textwidth}{@{}l >{\raggedright\arraybackslash}X >{\raggedright\arraybackslash}X@{}}
\toprule
Round & Data added & Main method added \\
\midrule
v1 & RXTE Standard-2 spectra, 8 sources & pipeline, LOSO, LR/RF on representations A and B \\
v2 & 31 sources, 456 observations; Standard-1 dead time & two-colour baseline H, source bootstrap, XSPEC check \\
v3 & 3--25\,keV grid; 14 BH candidates & nested increment tests (H+B vs H), source-level AUC \\
v4 & HEXTE; earlier gain epochs (v4b2); all 7949 epoch-5 pointings (v4b1) & 8-algorithm nested search, burst detector, source-equal weights \\
v5--v6 & Standard-1 timing features; 23 external sources & cospectral rms $T_1$--$T_3$; all-pairs estimator; $\nu_c$; permutation test; CCTLR; conformal sets \\
v7 & MINBAR; Cyg X-1, LMC X-1, LMC X-3; MAXI & rms state classes; burst matching; MAXI pipeline and reproduction of de Beurs et al. \\
v8 & GS 1354$-$64, SS 433; RXTE event mode; HI4PI & state-stratified bootstrap; high-frequency cospectra; $N_\mathrm{H}$ correction \\
v9--v10 & NICER N1 (prefix, then full files) & NICER streaming parser and features; pooled RXTE+NICER tests \\
v11 & NICER N2 (unseen observations) & frozen out-of-source prediction; Lorentzian fits \\
v12 & 3C50 background for N1+N2 & background correction and screening \\
v13 & NICER N3 (unseen observations) & everything frozen before download \\
v14 & -- & cross-round summary only \\
\bottomrule
\end{tabularx}
\end{table}

\subsection{Repository layout and conventions}\label{sec:layout}
\begin{itemize}
\item \textbf{Version switch.} Every script sets or reads the environment variable \code{XRB_VERSION}, and
\file{scripts/config.py} activates the corresponding parameter block. Later versions extend earlier blocks, and each new
block is checked not to change the values of earlier versions. \file{scripts/common.py} maps the version to the output
folders: \file{data/<v>/}, \file{results/<v>/}, \file{figures/<v>/} and \file{logs/<v>/} (v1 uses the top-level folders).
\item \textbf{Raw data.} Raw data stay out of git. RXTE products of v1--v3 and the v8 event files are in
\file{data/raw/}. Large sets live under \code{XRB_EXTERNAL_RAW} (default \file{~/xrb-pilot-data/raw}): v4b1, the NICER
prefixes and tails, the v13 cleaned files and the 3C50 spectra (\file{~/xrb-pilot-data/nicer_bkg_spectra}). The NICER
CALDB is under \file{~/xrb-pilot-data/caldb}.
\item \textbf{In git.} Everything derived is committed: source and candidate tables, every tried observation with its
status and exclusion reason, feature tables, out-of-fold predictions, result tables, figures and logs.
\item \textbf{Randomness.} The global seed is \code{SEED = 42}. Sampling seeds are offsets of it (\cref{tab:seeds}); the
bootstrap, permutations and models use it directly.
\item \textbf{Shared modules.} \file{common.py} (download and progress log), \file{spectra.py} (Standard-2 reader without
HEASoft), \file{config.py} (settings) and the version libraries \file{v3lib.py}--\file{v11lib.py}. Each version library
is tested by a \file{test_vNlib.py}.
\end{itemize}
